\documentclass[11pt,a4paper]{article}
\usepackage[utf8]{inputenc}
\usepackage[T1]{fontenc}
\usepackage[french]{babel}
\usepackage{amsmath,amssymb,amsthm}
\usepackage{booktabs}
\usepackage{geometry}
\usepackage{hyperref}
\geometry{margin=2.5cm}

\title{Methodologie de modeles neuro-symboliques interpretable:\\
HyConEx, HyperLogic et DR-HyperCF}
\author{Projet de recherche -- Redaction technique}
\date{\today}

\begin{document}
\maketitle

\begin{abstract}
Ce document presente, dans un format de type article scientifique, la methodologie de trois approches
complementaires pour l'apprentissage interpretable sur donnees tabulaires:
(i) HyConEx, qui combine classification et generation de contrefactuels;
(ii) HyperLogic, qui integre un hyperreseau dans un apprentissage de regles differentiable;
et (iii) une variante operationnelle du projet, DR-HyperCF, qui unifie extraction de regles et
generation de contrefactuels en espace binaire avec projection continue finale.
Nous detaillons pour chaque approche l'architecture, les objectifs d'optimisation, le pipeline de donnees,
ainsi que les protocoles d'evaluation.
\end{abstract}

\section{Introduction}
Les modeles interpretable doivent concilier deux exigences en tension:
la \emph{performance predictive} et la \emph{lisibilite des decisions}.
Dans les architectures de type neuro-symbolique, l'objectif est de conserver la puissance d'optimisation
des reseaux neuronaux tout en produisant des sorties explicables (regles, motifs, ou trajectoires
contrefactuelles).

Dans ce cadre, HyConEx, HyperLogic et DR-HyperCF peuvent etre vus comme trois niveaux de design:
\begin{itemize}
    \item HyConEx: architecture hybride classification+CF sur variables continues;
    \item HyperLogic: generation de poids d'un reseau de regles via hyperreseau (diversite des regles);
    \item DR-HyperCF (ce projet): version entierement binaire pendant l'apprentissage, orientee regles et CF.
\end{itemize}

\section{Methodologie HyConEx}
\subsection{Architecture}
HyConEx adopte trois composantes:
\begin{enumerate}
    \item un encodeur tabulaire $E(\cdot)$ qui projette $x$ dans un latent $z$;
    \item un hypermodule qui produit des parametres de classification dynamiques conditionnes par $z$;
    \item un generateur de contrefactuels conditionne par une classe cible $y_t$.
\end{enumerate}

Pour une entree $x \in \mathbb{R}^d$:
\begin{align}
z &= E(x),\\
\theta(x) &= H(z),\\
f(x) &= \text{Classifier}_{\theta(x)}(z),\\
x_{cf} &= \text{clip}\left(x + G\left([z,\text{onehot}(y_t)]\right), 0,1\right).
\end{align}

\subsection{Objectif d'entrainement}
La perte combine performance de classification et qualite contrefactuelle:
\begin{equation}
\mathcal{L}_{HyConEx}
= \underbrace{\text{CE}(f(x),y)}_{\text{classification}}
+ \lambda_{cf}\underbrace{\text{CE}(f(x_{cf}),y_t)}_{\text{validite CF}}
+ \lambda_1\|x_{cf}-x\|_1
+ \lambda_2\|x_{cf}-x\|_2^2.
\end{equation}

Le compromis est pilote par $\lambda_{cf}, \lambda_1, \lambda_2$.
Une normalisation des variables dans $[0,1]$ est generalement necessaire.

\subsection{Evaluation}
HyConEx est evalue selon deux axes:
\begin{itemize}
    \item \textbf{Classification}: accuracy, AUROC, matrice de confusion;
    \item \textbf{Contrefactuels}: validite ($\hat{y}_{cf}=y_t$), proximite ($L_1$), parcimonie (nombre de variables modifiees).
\end{itemize}

\section{Methodologie HyperLogic}
\subsection{Main network de regles}
HyperLogic part d'un reseau de regles differentiable (souvent DR-Net like), dans lequel chaque neurone
cache encode une regle candidate et la couche de sortie combine ces regles.
L'entree est typiquement binaire/discretisee.

\subsection{Hyperreseau generateur de poids}
Au lieu d'apprendre un unique jeu de poids $\theta$, HyperLogic apprend une distribution de poids:
\begin{equation}
\theta = G(\epsilon), \quad \epsilon \sim p(\epsilon),
\end{equation}
avec aggregation Monte-Carlo de plusieurs experts:
\begin{equation}
f_M(x)=\frac{1}{M}\sum_{m=1}^{M} f(x;\theta^{(m)}).
\end{equation}

Cette strategie augmente la flexibilite et la diversite des jeux de regles extraits.

\subsection{Perte et regularisation}
La fonction de cout combine:
\begin{itemize}
    \item une perte de tache ($\mathcal{L}_{task}$, ex. BCE/CE),
    \item une regularisation de diversite (souvent entropique/KL),
    \item une regularisation de sparsite des regles (souvent norme $\ell_1$).
\end{itemize}

Forme generale:
\begin{equation}
\mathcal{L}_{HyperLogic}
= \mathcal{L}_{task}
+ \lambda_{div}\mathcal{R}_{div}
+ \lambda_{sp}\mathcal{R}_{sp}.
\end{equation}

\subsection{Exploitation}
HyperLogic permet:
\begin{enumerate}
    \item d'echantillonner de nombreux ensembles de regles;
    \item de selectionner un ensemble optimal selon un critere (accuracy/complexite);
    \item d'eventuellement faire un ensemble de regles (voting) pour renforcer la robustesse.
\end{enumerate}

\section{Methodologie DR-HyperCF (nouveau module)}
\subsection{Objectif de design}
Le module DR-HyperCF introduit dans le projet suit les contraintes suivantes:
\begin{itemize}
    \item pas de bruit d'entree explicite dans l'usage courant;
    \item binarisation en debut de pipeline;
    \item apprentissage entierement en espace binaire;
    \item hyperreseau generant \textbf{100\%} des poids du main network et de la tete CF;
    \item conversion binaire$\rightarrow$continu seulement en sortie pour l'interpretation.
\end{itemize}

\subsection{Pipeline de donnees}
Soit $x_{cont}$ l'entree continue:
\begin{equation}
x_{bin} = B(x_{cont}),
\end{equation}
ou $B$ est un binariseur par discretisation en bins + one-hot par variable.

Le binariseur stocke les intervalles afin de reconstruire ensuite une representation continue:
\begin{equation}
x_{cf,cont}=B^{-1}(x_{cf,bin}).
\end{equation}

\subsection{Architecture a deux tetes}
\paragraph{Tete 1: main network de regles}
Le main network DR-like produit les logits de classe a partir de $x_{bin}$.
Les regles sont extraites depuis ses poids (litteraux ON/OFF sur les bits binaires).

\paragraph{Tete 2: generation contrefactuelle binaire}
La tete CF prend $x_{bin}$ et $y_t$, puis genere $x_{cf,bin}$.
La validite est verifiee en repassant $x_{cf,bin}$ dans la tete regles.

\paragraph{Hyperreseau}
L'hyperreseau recoit des statistiques de batch sur $x_{bin}$ (moyenne, ecart-type) et produit:
\begin{equation}
(\theta_{main},\theta_{cf}) = \mathcal{H}(\text{stats}(x_{bin})).
\end{equation}

Il n'existe pas de fusion $\alpha$ avec des poids natifs: les poids effectifs sont entierement generes.

\subsection{Fonction de cout}
La perte totale implementee suit:
\begin{equation}
\mathcal{L}_{DR-HyperCF}
= \underbrace{\text{CE}(f_{main}(x_{bin}),y)}_{\mathcal{L}_{cls}}
+ \lambda_{cf}\underbrace{\text{CE}(f_{main}(x_{cf,bin}),y_t)}_{\mathcal{L}_{cf\_valid}}
+ \lambda_{flip}\underbrace{\|x_{cf,bin}-x_{bin}\|_1}_{\mathcal{L}_{flip}}
+ \lambda_{rule}\underbrace{\mathcal{R}_{rule}}_{\text{sparsite regles}}.
\end{equation}

\subsection{Interpretation des sorties}
\begin{itemize}
    \item \textbf{Prediction standard}: toujours issue de la tete regles.
    \item \textbf{Sortie CF}: proposition d'une modification binaire puis conversion continue pour usage humain.
    \item \textbf{Regles}: chaque \texttt{rule\_id} correspond a un neurone regle; les regles peuvent etre affichees en version compacte (top literals) ou complete (tous literals au-dessus d'un seuil).
\end{itemize}

\section{Comparaison synthetique des trois approches}
\begin{table}[h]
\centering
\begin{tabular}{@{}llll@{}}
\toprule
Critere & HyConEx & HyperLogic & DR-HyperCF \\
\midrule
Espace principal & Continu & Binaire/Discretise & Binaire \\
Hyperreseau & Oui (dynamique) & Oui (distribution de poids) & Oui (poids 100\%) \\
Contrefactuels & Oui (continu) & Pas nativement central & Oui (binaire puis continu) \\
Extraction de regles & Indirecte/moderee & Centrale & Centrale \\
Diversite de jeux de regles & Limitee & Forte (sampling experts) & Moderee (selon hypernet) \\
\bottomrule
\end{tabular}
\caption{Positionnement methodologique des trois modeles.}
\end{table}

\section{Discussion pratique}
Pour un usage appliqué:
\begin{enumerate}
    \item HyConEx est adapte lorsque la priorite est la qualite des CF continus et la performance predictive.
    \item HyperLogic est pertinent pour generer plusieurs ensembles de regles et analyser leur diversite.
    \item DR-HyperCF est un compromis operationnel du projet, focalise sur la lisibilite logique en binaire avec projection continue finale.
\end{enumerate}

\section{Conclusion}
Les trois methodes s'inscrivent dans une meme logique:
integrer l'optimisation differentiable avec des structures interpretables.
HyConEx met l'accent sur le couple classification--contrefactuels continus,
HyperLogic sur la diversite de regles via hyperreseau,
et DR-HyperCF sur une chaine entierement binaire pour faciliter l'extraction de regles
et l'explicabilite des modifications contrefactuelles.

\end{document}
